\subsection{Instruction Header Formats}

Instruction framing is byte-oriented. Byte 0 identifies extrashort and short instructions. Extended
instructions use byte 0 and byte 1: byte 0 supplies the framing and encoded length, while the remaining
bits of those bytes carry the first fixed bits that select the (:term:base.terms.opcode:) and the first
operand-selector (:term:base.terms.field|plural:).

For an extended instruction, byte0[5:2] is (:instruction-header-field-target:base.encoding.instruction.EXTENDED_HEADER.L:)\texttt{L} and the (:term:base.terms.encoded_instruction_length:) is $3+L$ bytes. The
(:term:base.terms.opcode_space:) contains fixed bits that select the (:term:base.terms.opcode:) and operand-selector (:term:base.terms.field|plural:);
it occupies three bytes for medium instructions, four bytes for long instructions, five bytes for extralong
instructions, and six bytes for xxlong instructions. An extended instruction is invalid when its encoded
length is shorter than the (:term:base.terms.opcode_space:) selected by its (:term:base.terms.header:).

\texttt{LEN n, <instruction>} requests an encoded extended instruction length of \texttt{n} bytes, where
\texttt{3 <= n <= 18}. The requested length must cover the (:term:base.terms.required_instruction_length:). Without \texttt{LEN}, the
assembler emits that required length. It fills requested trailing (:term:base.terms.padding:) with zero. A disassembler emits
\texttt{LEN n,} when the (:term:base.terms.encoded_instruction_length:) exceeds the (:term:base.terms.required_instruction_length:) so reassembly preserves
the length; (:term:base.terms.padding:) byte values are not (:term:base.terms.preserved:). Extrashort and short encodings retain their fixed lengths.

\begin{BedrockListedBitDiagram}{Instruction Start Byte Formats}[fig:instruction-header-byte-formats]
\BedrockBitRow{extrashort byte}{%
\BedrockBitFixed{0}{1}
\BedrockBitFieldText{enc.}{7}
}
\BedrockBitFieldRow{short bytes}{\BedrockByteRowLabels{0}{2}}{%
\BedrockBitFixed{10}{2}
\BedrockBitFieldText{enc.}{6}
\BedrockBitGap{1}
\BedrockBitFieldText{enc.}{8}
}
\BedrockBitFieldRow{extended start bytes}{\BedrockByteRowLabels{0}{2}}{%
\BedrockBitFixed{11}{2}
\BedrockBitFieldText{L}{4}
\BedrockBitFieldText{}{2}
\BedrockBitGap{1}
\BedrockBitFieldText{enc.}{8}
}
\end{BedrockListedBitDiagram}

Figure~\ref{fig:instruction-header-byte-formats} is read from high to low bit. \texttt{enc.} denotes
instruction-specific bits within the (:term:base.terms.opcode_space:): fixed (:term:base.terms.opcode:) bits and operand-selector (:term:base.terms.field|plural:). In byte 0, bit 7
selects extrashort framing; bits 7..6 select short (\texttt{10}) or extended (\texttt{11}) framing. For extended
instructions, (:ref:base.encoding.instruction.EXTENDED_HEADER.L:) directly encodes the (:term:base.terms.encoded_instruction_length:) as $3+L$ bytes.

The extended class selector is byte0[1:0] followed by byte1[7:4]. Values 0 through 59 select a medium
(:term:base.terms.opcode_space:), and values 60 through 62 select a long (:term:base.terms.opcode_space:). When the selector is 63,
byte1[3:2] completes an eight-bit class prefix: \texttt{1111110x} selects extralong and
\texttt{11111111} selects xxlong. The \texttt{11111110} prefix retains extralong framing but is
(:term:base.terms.reserved:) for custom extensions and is not assigned by the standard ISA.

\BedrockTableCaption{Encoded Instruction Length Truth Table}
\begin{BedrockLongTable}{@{}>{\raggedright\arraybackslash}p{1.25in}>{\raggedright\arraybackslash}p{1.25in}>{\raggedright\arraybackslash}p{0.55in}@{}}
\toprule
\rowcolor{BedrockHeaderFill}
\textbf{Byte 0 Pattern} & \textbf{Byte 1 Pattern} & \textbf{Bytes}\\
\midrule
\endfirsthead
\multicolumn{3}{l}{\scriptsize\itshape Table \theBedrockTable\ (continued)}\\
\toprule
\rowcolor{BedrockHeaderFill}
\textbf{Byte 0 Pattern} & \textbf{Byte 1 Pattern} & \textbf{Bytes}\\
\midrule
\endhead
\texttt{0xxxxxxx} & -- & 1\\
\texttt{10xxxxxx} & \texttt{xxxxxxxx} & 2\\
\texttt{110000xx} & \texttt{xxxxxxxx} & 3\\
\texttt{110001xx} & \texttt{xxxxxxxx} & 4\\
\texttt{110010xx} & \texttt{xxxxxxxx} & 5\\
\texttt{110011xx} & \texttt{xxxxxxxx} & 6\\
\texttt{110100xx} & \texttt{xxxxxxxx} & 7\\
\texttt{110101xx} & \texttt{xxxxxxxx} & 8\\
\texttt{110110xx} & \texttt{xxxxxxxx} & 9\\
\texttt{110111xx} & \texttt{xxxxxxxx} & 10\\
\texttt{111000xx} & \texttt{xxxxxxxx} & 11\\
\texttt{111001xx} & \texttt{xxxxxxxx} & 12\\
\texttt{111010xx} & \texttt{xxxxxxxx} & 13\\
\texttt{111011xx} & \texttt{xxxxxxxx} & 14\\
\texttt{111100xx} & \texttt{xxxxxxxx} & 15\\
\texttt{111101xx} & \texttt{xxxxxxxx} & 16\\
\texttt{111110xx} & \texttt{xxxxxxxx} & 17\\
\texttt{111111xx} & \texttt{xxxxxxxx} & 18\\
\bottomrule
\end{BedrockLongTable}

\BedrockTableCaption{Opcode-Space Class Summary}
\begin{BedrockTabular}{@{}>{\raggedright\arraybackslash}p{0.72in}>{\raggedright\arraybackslash}p{1.10in}>{\raggedright\arraybackslash}p{2.05in}>{\raggedright\arraybackslash}p{0.62in}>{\raggedright\arraybackslash}p{1.15in}@{}}
\toprule
\rowcolor{BedrockHeaderFill}
\textbf{Class} & \textbf{Opcode-Space Bytes} & \textbf{Header Selector} & \textbf{Minimum Bytes} & \textbf{Validity}\\
\midrule
\texttt{extrashort} & byte 0 & \texttt{0xxxxxxx} & 1 & exactly 1 byte\\
\texttt{short} & bytes 0--1 & \texttt{10xxxxxx} & 2 & exactly 2 bytes\\
\texttt{medium} & bytes 0--2 & extended selector 0--59 & 3 & encoded length at least 3 bytes\\
\texttt{long} & bytes 0--3 & extended selector 60--62 & 4 & encoded length at least 4 bytes\\
\texttt{extralong} & bytes 0--4 & extended selector 63, prefix \texttt{1111110x} & 5 & encoded length at least 5 bytes\\
\texttt{xxlong} & bytes 0--5 & all-one eight-bit prefix & 6 & encoded length at least 6 bytes\\
\bottomrule
\end{BedrockTabular}
